\section{Prediction and posteriors (base case)}

\subsection{The information at forecast time}
At the close of day $t$ the forecast may use $\mathcal G_t$: the market history (through the forward pass only), today's characteristics $x_{i,t}$, and the frozen parameters. The gate is the $\F_t$-measurable probability
\begin{equation}
\pi_k(t)=\mathbb P(z_{t+1}=k\mid\F_t)=\sum_jA_{jk}\,\xi_{t\mid t}(j).
\end{equation}

\subsection{The predictive distribution of one stock}
Apply the sum rule over $z_{t+1}$:
\begin{equation}\label{eq:pred}
p\big(r_{i,t+1}\mid x_{i,t},\F_t\big)=\sum_{k=1}^K\pi_k(t)\,\N\big(r_{i,t+1};\mu_k(x_{i,t}),s_k^2\big).
\end{equation}
This is a Gaussian mixture, so it can be skewed and fat-tailed even though every expert is Gaussian. It is the full forecast; everything below summarises it.

\subsection{The point forecast}
By the law of total expectation,
\begin{equation}
\hat r_{i,t+1}=\E\big[r_{i,t+1}\mid x_{i,t},\F_t\big]=\sum_k\pi_k(t)\,\mu_k(x_{i,t})=f_0(x_{i,t})+\sum_k\pi_k(t)\,c_k(x_{i,t}).
\end{equation}
The last step uses $\sum_k\pi_k(t)=1$: the base comes out of the sum, so the forecast is the base plus the gate-weighted corrections. The rank IC is computed from $\hat r_{i,t+1}$.

\subsection{The predictive variance}
The second moment is $\E[r^2\mid\cdot]=\sum_k\pi_k(s_k^2+\mu_k^2)$. Subtracting $\hat r^{\,2}$ and rearranging gives the law of total variance:
\begin{equation}
\Var\big(r_{i,t+1}\mid x_{i,t},\F_t\big)=\underbrace{\sum_k\pi_k(t)\,s_k^2}_{\text{noise within a regime}}+\underbrace{\sum_k\pi_k(t)\big(\mu_k(x_{i,t})-\hat r_{i,t+1}\big)^2}_{\text{experts disagree, regime uncertain}}.
\end{equation}
The first term remains even if tomorrow's regime were known. The second exists only because the regime is unknown and the experts predict different things; it vanishes if the gate is certain or if the experts agree for this stock.

\subsection{Covariance between stocks (per-date model)}
Two stocks share $z_{t+1}$ and are independent given it, so uncertainty about the regime makes them co-move:
\begin{equation}
\Cov\big(r_{i,t+1},r_{j,t+1}\mid\cdot\big)=\sum_k\pi_k\,\mu_k(x_{i,t})\,\mu_k(x_{j,t})-\hat r_{i,t+1}\hat r_{j,t+1}=\sum_k\pi_k(t)\big(\mu_k(x_{i,t})-\hat r_{i,t+1}\big)\big(\mu_k(x_{j,t})-\hat r_{j,t+1}\big).
\end{equation}
This covariance is a property of the model; it is not affected by training with the per-row objective (Q24), which only discards it from the estimation. Under the alternative per-row \emph{model} it would be zero. It matters only for portfolios built from the predictive distribution; ranking uses $\hat r$ alone.

\subsection{The study's horizon: $h=5$}
The study uses a daily gate, a daily cross-section and a 5-day target (Q21, decided 2026-10-05). The derivations in these notes are written for one step, with $r_{i,t+1}$; in the study that step is the 5-day window $(t,t+5]$, and $r_{i,t+1}$ stands for the market-neutral return over it. The regime can change inside the window, and the regime forecast $j$ days ahead is $\mathbb P(z_{t+j}=k\mid\F_t)=[\xi_{t\mid t}\T A^j]_k$. The likelihood keeps a single regime for the whole window, which is a reasonable approximation because 5 days is short against typical regime durations of 20--50 days (Section 5.3); the window-average weight below makes the point forecast exact even when the regime changes inside the window. The gate weight for the window is the average of the $j=1,\dots,5$ forecasts (decided 2026-10-05, Q27; Section~\ref{sec:predfilt}). Consecutive daily targets overlap by four days: statistics are overlap-corrected (Hansen--Hodrick), and the training likelihood is used for estimation, not inference.

\subsection{Every regime probability, and where it may be used}
\begin{center}
\small\renewcommand{\arraystretch}{1.3}
\begin{tabularx}{\textwidth}{@{}p{2.6cm}p{1.6cm}p{2.6cm}Xp{1.4cm}@{}}
\toprule
Quantity & Measurable w.r.t. & Computed from & Used for & At forecast time \\
\midrule
filtered $\xi_{t\mid t}$ & $\F_t$ & forward pass ($\hat\alpha_t$) & building the gate & yes \\
gate $\bar w_k(t)$ ($=\pi_k(t)$ if $h=1$) & $\F_t$ & $\xi_{t\mid t}$ and $A$ & weighting the experts & yes \\
smoothed $\gamma_t$ (gate) & $\F_T$ & $\hat\alpha_t\hat\beta_t$ & Baum--Welch M-step, training block & \textbf{no} \\
pairwise $\xi_t(j,k)$ & $\F_T$ & $\hat\alpha,A,\eta,\hat\beta$ & updating $A$, training block & \textbf{no} \\
responsibility $\gamma_{i,t}(k)$ & $\mathcal G_{t+1}$ & Bayes per stock-day, uses $r_{i,t+1}$ & weighting expert gradients in training & \textbf{no} \\
\bottomrule
\end{tabularx}
\end{center}

\subsection{Why smoothed probabilities leak the future}
\begin{equation}
\gamma_t(k)\propto\alpha_t(k)\,\beta_t(k),\qquad\beta_t(k)=p(m_{t+1},\dots,m_T\mid z_t=k).
\end{equation}
Suppose the market crashes on day $t+3$. That crash is far more likely under the stressed regime, so $\beta_t$ shifts weight to the stressed regime on day $t$, three days before the crash, when no past data pointed to it. A gate built on smoothed probabilities therefore enters crashes in advance: the backtest looks excellent and means nothing. In the language of Section 1.4, the smoothed probability is not $\F_t$-measurable, so a forecast built on it is not adapted.

\paragraph{Leaks that are easy to miss.}
\begin{itemize}
  \item Refitting Baum--Welch on data that include validation or test dates. The M-step uses smoothed posteriors, so future data enter $\theta_g$.
  \item Training on smoothed weights and testing on filtered ones. This is not look-ahead at test time, but the experts face a different kind of input than they were trained on. That is why the filtered series is used on training dates too.
  \item Using the responsibilities $\gamma$ as features at forecast time. They contain the target itself.
  \item The purge gap (Q22, open). The forward pass must be carried continuously through the dates removed between training and test.
\end{itemize}

\subsection{The gate weight for an $h$-day target (Q27, decided)}\label{sec:predfilt}
Three $\F_t$-measurable candidates are available at the close of day $t$, all built from the filtered probability:
\begin{itemize}
  \item the \emph{filtered} probability $\xi_{t\mid t}(k)$, today's regime (used by \path{Model_Derivation_BaseCase.md} before this decision, and by the code's Hamilton gate);
  \item the \emph{one-step predicted} probability $\pi_k(t)=[\xi_{t\mid t}\T A]_k$, tomorrow's regime (Section 6 with $h=1$);
  \item the \emph{window average}, the expected share of the window $(t,t+h]$ spent in regime $k$:
\end{itemize}
\begin{equation}\label{eq:wbar}
\bar w_k(t)=\frac1h\sum_{j=1}^{h}\big[\xi_{t\mid t}\T A^{\,j}\big]_k.
\end{equation}

\paragraph{Why the window average is exact for the point forecast.} The target is the sum of $h$ daily returns. Suppose each day's expected return is set by that day's regime, with regime $k$ contributing $\mu_k(x_{i,t})/h$ per day. By the law of total expectation, day by day,
\begin{equation}
\E\big[y_{i,t}\mid\F_t\big]=\sum_{j=1}^h\sum_k\mathbb P\big(z_{t+j}=k\mid\F_t\big)\frac{\mu_k(x_{i,t})}{h}=\sum_k\bar w_k(t)\,\mu_k(x_{i,t}).
\end{equation}
So the gate-weighted forecast with \eqref{eq:wbar} is exactly the expected $h$-day return. The expected within-regime noise also adds up over days and is matched in the same way. The filtered weight assumes the regime never changes during the window, and the one-step weight assumes it changes at most once. For $h=1$, \eqref{eq:wbar} reduces to the one-step predicted probability, so every one-step formula in these notes is the $h=1$ case of this rule.

\paragraph{How large the differences are.} With the example $A$ of Section 5.3 and a filtered belief of certain stress, the weight on stress is $1.00$ (filtered), $0.95$ (one-step) and $0.86$ (window average with $h=5$). A descriptive fit of a two-state Hamilton model to the CRSP S\&P~500 index, 2015--2024, gives an average gap of $0.055$ between the filtered and window-average weights on the stressed regime, with a maximum of $0.10$. The gap is largest when the filter is sure of stress, because stress tends to end.

\paragraph{Decision.} The window average is used on training and test dates alike (Q27, decided 2026-10-05). Costs accepted: slightly softer weights, and dependence on a well-estimated $A$, whose errors compound over $h$ steps. The likelihood with these weights remains the one-regime-per-window approximation: the exact distribution of an $h$-day return is a mixture over regime \emph{paths}.
\newpage
